% --- Archon physics formalization source begin ---
% archon:physics
% archon:covers PhyXMiniProblems/problem_phyx_mini_0755.lean
% archon:source-report reports/phyx_mini/problem_phyx_mini_0755.source.json

\chapter{Physics problem phyx\_mini\_0755}
\label{ch:PhyXMiniProblems_problem_phyx_mini_0755}

\paragraph{Problem source.}
\#\# Physical scenario

Three astronauts, propelled by jet backpacks, push and guide a \$120\textbackslash{} kg\$ asteroid toward a processing dock, exerting the forces shown in figure, with \$F\_1 = 32\textbackslash{} N\$, \$F\_2 = 55\textbackslash{} N\$, \$F\_3 = 41\textbackslash{} N\$, \$\textbackslash{}theta\_1 = 30\textasciicircum{}\textbackslash{}circ\$, and \$\textbackslash{}theta\_3 = 60\textasciicircum{}\textbackslash{}circ\$.

\#\# Question

What is the asteroid's acceleration as a magnitude relative to the positive direction of the \$x\$ axis?

\#\# Answer choices

- A: A: \$0.76\textbackslash{} m/s\textasciicircum{}2\$

- B: B: \$0.82\textbackslash{} m/s\textasciicircum{}2\$

- C: C: \$0.88\textbackslash{} m/s\textasciicircum{}2\$

- D: D: \$0.92\textbackslash{} m/s\textasciicircum{}2\$

\#\# Figure caption (auxiliary; use the image as primary evidence)

The image depicts three astronauts pushing a large, irregularly shaped object in space. Each astronaut is emitting thruster exhaust. The scene has an overlaid x-y coordinate system.

There are three force vectors labeled \textbackslash{}(\textbackslash{}vec\{F\}\_1\textbackslash{}), \textbackslash{}(\textbackslash{}vec\{F\}\_2\textbackslash{}), and \textbackslash{}(\textbackslash{}vec\{F\}\_3\textbackslash{}), each acting in different directions. The angles \textbackslash{}(\textbackslash{}theta\_1\textbackslash{}) and \textbackslash{}(\textbackslash{}theta\_3\textbackslash{}) are shown between the vectors and the y-axis.

The force vectors and angles suggest a study of dynamics, focusing on the forces applied by the astronauts on the object.

\#\# Dataset metadata

Dynamics; Spatial Relation Reasoning, Multi-Formula Reasoning

\paragraph{Recorded answer/context.}
C: C: \$0.88\textbackslash{} m/s\textasciicircum{}2\$

\paragraph{Figure/image path.}
/root/proposal\_for\_physic/science-mango/phyx\_mini\_run/phyx\_data/test\_image/755.png

\paragraph{Formalization target.}
create a compiling Lean file with sorry bodies at `PhyXMiniProblems/problem\_phyx\_mini\_0755.lean`. The Lean declarations must preserve the physical quantities, dimensions or dimensional roles, figure labels, governing-law hypotheses, and final relation expressed by this problem.
Use Mathlib/Physlib names found through LeanExplore where available. If a physics API is missing, introduce faithful local abstractions rather than scalar placeholder aliases.

\begin{theorem}[Physics formalization target]
\label{thm:physics:phyx_mini_0755:target}
\lean{PhyXMiniProblems.ProblemPhyXMini0755.problem_phyx_mini_0755}
\uses{def:physics:phyx-mini-0755:phyxminiproblems-problemphyxmini0755-accelerationmagnitudeinmeterspersecondsquared, def:physics:phyx-mini-0755:phyxminiproblems-problemphyxmini0755-asteroidpushsetup, def:physics:phyx-mini-0755:phyxminiproblems-problemphyxmini0755-matchesasteroidguidancescenario, def:physics:phyx-mini-0755:phyxminiproblems-problemphyxmini0755-matchesproblemreadouts, def:physics:phyx-mini-0755:phyxminiproblems-problemphyxmini0755-matchesprimaryfigure, def:physics:phyx-mini-0755:phyxminiproblems-problemphyxmini0755-matchesfigureforcegeometry, def:physics:phyx-mini-0755:phyxminiproblems-problemphyxmini0755-satisfiesnewtoniandynamics, def:physics:phyx-mini-0755:phyxminiproblems-problemphyxmini0755-exactaccelerationmagnitudeinmeterspersecondsquared, def:physics:phyx-mini-0755:phyxminiproblems-problemphyxmini0755-matchesdisplayedacceleration, def:physics:phyx-mini-0755:phyxminiproblems-problemphyxmini0755-isuniqueclosestaccelerationchoice}
The assigned autoformalize agent should translate this physics problem into Lean declarations in the covered file, with theorem and lemma proof bodies written as `by sorry`.
\end{theorem}
\begin{proof}
This is an autoformalization task, not a proof task. Produce statements that can later be proved by the physics prover without weakening the physical model.
\end{proof}

% --- Archon declaration topology begin ---
\section{Declaration topology}
\begin{definition}[acceleration Dimension]
  \label{def:physics:phyx-mini-0755:phyxminiproblems-problemphyxmini0755-accelerationdimension}
  \lean{PhyXMiniProblems.ProblemPhyXMini0755.accelerationDimension}
  The physical dimension of acceleration, L T⁻²
\end{definition}
\begin{proof}
  This is the declared model definition of acceleration Dimension.
\end{proof}
\begin{definition}[force Dimension]
  \label{def:physics:phyx-mini-0755:phyxminiproblems-problemphyxmini0755-forcedimension}
  \lean{PhyXMiniProblems.ProblemPhyXMini0755.forceDimension}
  \uses{def:physics:phyx-mini-0755:phyxminiproblems-problemphyxmini0755-accelerationdimension}
  The physical dimension of force, M L T⁻²
\end{definition}
\begin{proof}
  This is the declared model definition of force Dimension.
\end{proof}
\begin{definition}[Mass Quantity]
  \label{def:physics:phyx-mini-0755:phyxminiproblems-problemphyxmini0755-massquantity}
  \lean{PhyXMiniProblems.ProblemPhyXMini0755.MassQuantity}
  A unit-independent nonnegative physical mass
\end{definition}
\begin{proof}
  This is the declared model definition of Mass Quantity.
\end{proof}
\begin{definition}[Planar Force Quantity]
  \label{def:physics:phyx-mini-0755:phyxminiproblems-problemphyxmini0755-planarforcequantity}
  \lean{PhyXMiniProblems.ProblemPhyXMini0755.PlanarForceQuantity}
  \uses{def:physics:phyx-mini-0755:phyxminiproblems-problemphyxmini0755-forcedimension}
  A unit-independent force vector in the plane of the supplied figure
\end{definition}
\begin{proof}
  This is the declared model definition of Planar Force Quantity.
\end{proof}
\begin{definition}[Planar Acceleration Quantity]
  \label{def:physics:phyx-mini-0755:phyxminiproblems-problemphyxmini0755-planaraccelerationquantity}
  \lean{PhyXMiniProblems.ProblemPhyXMini0755.PlanarAccelerationQuantity}
  \uses{def:physics:phyx-mini-0755:phyxminiproblems-problemphyxmini0755-accelerationdimension}
  A unit-independent acceleration vector in the plane of the figure
\end{definition}
\begin{proof}
  This is the declared model definition of Planar Acceleration Quantity.
\end{proof}
\begin{definition}[x Axis]
  \label{def:physics:phyx-mini-0755:phyxminiproblems-problemphyxmini0755-xaxis}
  \lean{PhyXMiniProblems.ProblemPhyXMini0755.xAxis}
  Coordinate 0, the positive horizontal axis pointing right
\end{definition}
\begin{proof}
  This is the declared model definition of x Axis.
\end{proof}
\begin{definition}[y Axis]
  \label{def:physics:phyx-mini-0755:phyxminiproblems-problemphyxmini0755-yaxis}
  \lean{PhyXMiniProblems.ProblemPhyXMini0755.yAxis}
  Coordinate 1, the positive vertical axis pointing up
\end{definition}
\begin{proof}
  This is the declared model definition of y Axis.
\end{proof}
\begin{definition}[mass In Kilograms]
  \label{def:physics:phyx-mini-0755:phyxminiproblems-problemphyxmini0755-massinkilograms}
  \lean{PhyXMiniProblems.ProblemPhyXMini0755.massInKilograms}
  \uses{def:physics:phyx-mini-0755:phyxminiproblems-problemphyxmini0755-massquantity}
  Coherent-SI kilogram readout of a physical mass
\end{definition}
\begin{proof}
  This is the declared model definition of mass In Kilograms.
\end{proof}
\begin{definition}[force Vector In Newtons]
  \label{def:physics:phyx-mini-0755:phyxminiproblems-problemphyxmini0755-forcevectorinnewtons}
  \lean{PhyXMiniProblems.ProblemPhyXMini0755.forceVectorInNewtons}
  \uses{def:physics:phyx-mini-0755:phyxminiproblems-problemphyxmini0755-planarforcequantity}
  Coherent-SI planar-force readout, whose coordinates are in newtons
\end{definition}
\begin{proof}
  This is the declared model definition of force Vector In Newtons.
\end{proof}
\begin{definition}[force Magnitude In Newtons]
  \label{def:physics:phyx-mini-0755:phyxminiproblems-problemphyxmini0755-forcemagnitudeinnewtons}
  \lean{PhyXMiniProblems.ProblemPhyXMini0755.forceMagnitudeInNewtons}
  \uses{def:physics:phyx-mini-0755:phyxminiproblems-problemphyxmini0755-planarforcequantity, def:physics:phyx-mini-0755:phyxminiproblems-problemphyxmini0755-forcevectorinnewtons}
  Euclidean magnitude in newtons of a physical planar force
\end{definition}
\begin{proof}
  This is the declared model definition of force Magnitude In Newtons.
\end{proof}
\begin{definition}[acceleration Vector In Meters Per Second Squared]
  \label{def:physics:phyx-mini-0755:phyxminiproblems-problemphyxmini0755-accelerationvectorinmeterspersecondsquared}
  \lean{PhyXMiniProblems.ProblemPhyXMini0755.accelerationVectorInMetersPerSecondSquared}
  \uses{def:physics:phyx-mini-0755:phyxminiproblems-problemphyxmini0755-planaraccelerationquantity}
  Coherent-SI acceleration vector, in metres per second squared
\end{definition}
\begin{proof}
  This is the declared model definition of acceleration Vector In Meters Per Second Squared.
\end{proof}
\begin{definition}[acceleration Magnitude In Meters Per Second Squared]
  \label{def:physics:phyx-mini-0755:phyxminiproblems-problemphyxmini0755-accelerationmagnitudeinmeterspersecondsquared}
  \lean{PhyXMiniProblems.ProblemPhyXMini0755.accelerationMagnitudeInMetersPerSecondSquared}
  \uses{def:physics:phyx-mini-0755:phyxminiproblems-problemphyxmini0755-planaraccelerationquantity, def:physics:phyx-mini-0755:phyxminiproblems-problemphyxmini0755-accelerationvectorinmeterspersecondsquared}
  Euclidean acceleration magnitude in metres per second squared
\end{definition}
\begin{proof}
  This is the declared model definition of acceleration Magnitude In Meters Per Second Squared.
\end{proof}
\begin{definition}[Astronaut Label]
  \label{def:physics:phyx-mini-0755:phyxminiproblems-problemphyxmini0755-astronautlabel}
  \lean{PhyXMiniProblems.ProblemPhyXMini0755.AstronautLabel}
  The three astronauts shown pushing the asteroid
\end{definition}
\begin{proof}
  This is the declared model definition of Astronaut Label.
\end{proof}
\begin{definition}[Applied Force Label]
  \label{def:physics:phyx-mini-0755:phyxminiproblems-problemphyxmini0755-appliedforcelabel}
  \lean{PhyXMiniProblems.ProblemPhyXMini0755.AppliedForceLabel}
  The three force-vector labels printed in the primary image
\end{definition}
\begin{proof}
  This is the declared model definition of Applied Force Label.
\end{proof}
\begin{definition}[force Exerted By]
  \label{def:physics:phyx-mini-0755:phyxminiproblems-problemphyxmini0755-forceexertedby}
  \lean{PhyXMiniProblems.ProblemPhyXMini0755.forceExertedBy}
  \uses{def:physics:phyx-mini-0755:phyxminiproblems-problemphyxmini0755-astronautlabel, def:physics:phyx-mini-0755:phyxminiproblems-problemphyxmini0755-appliedforcelabel}
  The force label associated with each astronaut's push
\end{definition}
\begin{proof}
  This is the declared model definition of force Exerted By.
\end{proof}
\begin{definition}[Figure Angle Label]
  \label{def:physics:phyx-mini-0755:phyxminiproblems-problemphyxmini0755-figureanglelabel}
  \lean{PhyXMiniProblems.ProblemPhyXMini0755.FigureAngleLabel}
  The two named angles printed between the oblique arrows and positive x
\end{definition}
\begin{proof}
  This is the declared model definition of Figure Angle Label.
\end{proof}
\begin{definition}[Diagram Axis]
  \label{def:physics:phyx-mini-0755:phyxminiproblems-problemphyxmini0755-diagramaxis}
  \lean{PhyXMiniProblems.ProblemPhyXMini0755.DiagramAxis}
  Coordinate-axis labels printed over the asteroid
\end{definition}
\begin{proof}
  This is the declared model definition of Diagram Axis.
\end{proof}
\begin{definition}[Arrow Direction]
  \label{def:physics:phyx-mini-0755:phyxminiproblems-problemphyxmini0755-arrowdirection}
  \lean{PhyXMiniProblems.ProblemPhyXMini0755.ArrowDirection}
  Qualitative directions needed to describe the axes and force arrows
\end{definition}
\begin{proof}
  This is the declared model definition of Arrow Direction.
\end{proof}
\begin{definition}[Rotation Sense]
  \label{def:physics:phyx-mini-0755:phyxminiproblems-problemphyxmini0755-rotationsense}
  \lean{PhyXMiniProblems.ProblemPhyXMini0755.RotationSense}
  Sense in which an angle is measured from the positive x-axis
\end{definition}
\begin{proof}
  This is the declared model definition of Rotation Sense.
\end{proof}
\begin{definition}[expected Force Direction]
  \label{def:physics:phyx-mini-0755:phyxminiproblems-problemphyxmini0755-expectedforcedirection}
  \lean{PhyXMiniProblems.ProblemPhyXMini0755.expectedForceDirection}
  \uses{def:physics:phyx-mini-0755:phyxminiproblems-problemphyxmini0755-appliedforcelabel, def:physics:phyx-mini-0755:phyxminiproblems-problemphyxmini0755-arrowdirection}
  The qualitative force direction visibly associated with each arrow
\end{definition}
\begin{proof}
  This is the declared model definition of expected Force Direction.
\end{proof}
\begin{definition}[force For Angle]
  \label{def:physics:phyx-mini-0755:phyxminiproblems-problemphyxmini0755-forceforangle}
  \lean{PhyXMiniProblems.ProblemPhyXMini0755.forceForAngle}
  \uses{def:physics:phyx-mini-0755:phyxminiproblems-problemphyxmini0755-appliedforcelabel, def:physics:phyx-mini-0755:phyxminiproblems-problemphyxmini0755-figureanglelabel}
  The force arrow whose inclination is denoted by each angle label
\end{definition}
\begin{proof}
  This is the declared model definition of force For Angle.
\end{proof}
\begin{definition}[expected Rotation Sense]
  \label{def:physics:phyx-mini-0755:phyxminiproblems-problemphyxmini0755-expectedrotationsense}
  \lean{PhyXMiniProblems.ProblemPhyXMini0755.expectedRotationSense}
  \uses{def:physics:phyx-mini-0755:phyxminiproblems-problemphyxmini0755-figureanglelabel, def:physics:phyx-mini-0755:phyxminiproblems-problemphyxmini0755-rotationsense}
  The rotation sense visibly associated with each angle label
\end{definition}
\begin{proof}
  This is the declared model definition of expected Rotation Sense.
\end{proof}
\begin{definition}[Asteroid Push Figure]
  \label{def:physics:phyx-mini-0755:phyxminiproblems-problemphyxmini0755-asteroidpushfigure}
  \lean{PhyXMiniProblems.ProblemPhyXMini0755.AsteroidPushFigure}
  \uses{def:physics:phyx-mini-0755:phyxminiproblems-problemphyxmini0755-astronautlabel, def:physics:phyx-mini-0755:phyxminiproblems-problemphyxmini0755-appliedforcelabel, def:physics:phyx-mini-0755:phyxminiproblems-problemphyxmini0755-figureanglelabel, def:physics:phyx-mini-0755:phyxminiproblems-problemphyxmini0755-diagramaxis, def:physics:phyx-mini-0755:phyxminiproblems-problemphyxmini0755-arrowdirection, def:physics:phyx-mini-0755:phyxminiproblems-problemphyxmini0755-rotationsense}
  Literal and qualitative content transcribed from image 755.png
\end{definition}
\begin{proof}
  This is the declared model definition of Asteroid Push Figure.
\end{proof}
\begin{definition}[Asteroid Push Setup]
  \label{def:physics:phyx-mini-0755:phyxminiproblems-problemphyxmini0755-asteroidpushsetup}
  \lean{PhyXMiniProblems.ProblemPhyXMini0755.AsteroidPushSetup}
  \uses{def:physics:phyx-mini-0755:phyxminiproblems-problemphyxmini0755-massquantity, def:physics:phyx-mini-0755:phyxminiproblems-problemphyxmini0755-planarforcequantity, def:physics:phyx-mini-0755:phyxminiproblems-problemphyxmini0755-planaraccelerationquantity, def:physics:phyx-mini-0755:phyxminiproblems-problemphyxmini0755-astronautlabel, def:physics:phyx-mini-0755:phyxminiproblems-problemphyxmini0755-appliedforcelabel, def:physics:phyx-mini-0755:phyxminiproblems-problemphyxmini0755-asteroidpushfigure}
  Independent physical fields of the experiment. In particular, acceleration and resultant force are not defined from any answer value; the hypotheses below constrain them through the physical laws
\end{definition}
\begin{proof}
  This is the declared model definition of Asteroid Push Setup.
\end{proof}
\begin{definition}[Matches Asteroid Guidance Scenario]
  \label{def:physics:phyx-mini-0755:phyxminiproblems-problemphyxmini0755-matchesasteroidguidancescenario}
  \lean{PhyXMiniProblems.ProblemPhyXMini0755.MatchesAsteroidGuidanceScenario}
  \uses{def:physics:phyx-mini-0755:phyxminiproblems-problemphyxmini0755-asteroidpushsetup}
  Qualitative physical scenario stated in the prose
\end{definition}
\begin{proof}
  This is the declared model definition of Matches Asteroid Guidance Scenario.
\end{proof}
\begin{definition}[Matches Problem Readouts]
  \label{def:physics:phyx-mini-0755:phyxminiproblems-problemphyxmini0755-matchesproblemreadouts}
  \lean{PhyXMiniProblems.ProblemPhyXMini0755.MatchesProblemReadouts}
  \uses{def:physics:phyx-mini-0755:phyxminiproblems-problemphyxmini0755-massinkilograms, def:physics:phyx-mini-0755:phyxminiproblems-problemphyxmini0755-forcemagnitudeinnewtons, def:physics:phyx-mini-0755:phyxminiproblems-problemphyxmini0755-asteroidpushsetup}
  Numerical readouts stated in the problem. Degrees have been converted to radians. No acceleration value or answer label occurs in this record
\end{definition}
\begin{proof}
  This is the declared model definition of Matches Problem Readouts.
\end{proof}
\begin{definition}[Matches Primary Figure]
  \label{def:physics:phyx-mini-0755:phyxminiproblems-problemphyxmini0755-matchesprimaryfigure}
  \lean{PhyXMiniProblems.ProblemPhyXMini0755.MatchesPrimaryFigure}
  \uses{def:physics:phyx-mini-0755:phyxminiproblems-problemphyxmini0755-expectedforcedirection, def:physics:phyx-mini-0755:phyxminiproblems-problemphyxmini0755-expectedrotationsense, def:physics:phyx-mini-0755:phyxminiproblems-problemphyxmini0755-asteroidpushsetup}
  Primary-image evidence. In particular, both labelled angles use positive x as their reference ray; this reading follows the raster rather than the auxiliary caption's incompatible mention of the y-axis
\end{definition}
\begin{proof}
  This is the declared model definition of Matches Primary Figure.
\end{proof}
\begin{definition}[Matches Figure Force Geometry]
  \label{def:physics:phyx-mini-0755:phyxminiproblems-problemphyxmini0755-matchesfigureforcegeometry}
  \lean{PhyXMiniProblems.ProblemPhyXMini0755.MatchesFigureForceGeometry}
  \uses{def:physics:phyx-mini-0755:phyxminiproblems-problemphyxmini0755-xaxis, def:physics:phyx-mini-0755:phyxminiproblems-problemphyxmini0755-yaxis, def:physics:phyx-mini-0755:phyxminiproblems-problemphyxmini0755-forcevectorinnewtons, def:physics:phyx-mini-0755:phyxminiproblems-problemphyxmini0755-forcemagnitudeinnewtons, def:physics:phyx-mini-0755:phyxminiproblems-problemphyxmini0755-asteroidpushsetup}
  Cartesian realization of the primary-image arrow geometry. The negative y component of F₃ expresses its clockwise inclination below positive x. These are general component laws in terms of the independently supplied magnitudes and angles, not acceleration conclusions
\end{definition}
\begin{proof}
  This is the declared model definition of Matches Figure Force Geometry.
\end{proof}
\begin{definition}[Satisfies Newtonian Dynamics]
  \label{def:physics:phyx-mini-0755:phyxminiproblems-problemphyxmini0755-satisfiesnewtoniandynamics}
  \lean{PhyXMiniProblems.ProblemPhyXMini0755.SatisfiesNewtonianDynamics}
  \uses{def:physics:phyx-mini-0755:phyxminiproblems-problemphyxmini0755-massinkilograms, def:physics:phyx-mini-0755:phyxminiproblems-problemphyxmini0755-forcevectorinnewtons, def:physics:phyx-mini-0755:phyxminiproblems-problemphyxmini0755-accelerationvectorinmeterspersecondsquared, def:physics:phyx-mini-0755:phyxminiproblems-problemphyxmini0755-asteroidpushsetup}
  Governing Newtonian dynamics: the resultant is the vector sum of the three applied forces, and the net force equals mass times acceleration in coherent SI components. No requested numerical acceleration occurs in these laws
\end{definition}
\begin{proof}
  This is the declared model definition of Satisfies Newtonian Dynamics.
\end{proof}
\begin{definition}[exact Acceleration Magnitude In Meters Per Second Squared]
  \label{def:physics:phyx-mini-0755:phyxminiproblems-problemphyxmini0755-exactaccelerationmagnitudeinmeterspersecondsquared}
  \lean{PhyXMiniProblems.ProblemPhyXMini0755.exactAccelerationMagnitudeInMetersPerSecondSquared}
  Closed unrounded acceleration magnitude obtained from the stated force components and 120 kg mass. This expression contains no setup acceleration field and therefore does not define the requested physical observable into the assumptions
\end{definition}
\begin{proof}
  This is the declared model definition of exact Acceleration Magnitude In Meters Per Second Squared.
\end{proof}
\begin{definition}[Answer Choice]
  \label{def:physics:phyx-mini-0755:phyxminiproblems-problemphyxmini0755-answerchoice}
  \lean{PhyXMiniProblems.ProblemPhyXMini0755.AnswerChoice}
  The four acceleration-answer labels displayed with the problem
\end{definition}
\begin{proof}
  This is the declared model definition of Answer Choice.
\end{proof}
\begin{definition}[acceleration Meters Per Second Squared]
  \label{def:physics:phyx-mini-0755:phyxminiproblems-problemphyxmini0755-answerchoice-accelerationmeterspersecondsquared}
  \lean{PhyXMiniProblems.ProblemPhyXMini0755.AnswerChoice.accelerationMetersPerSecondSquared}
  \uses{def:physics:phyx-mini-0755:phyxminiproblems-problemphyxmini0755-answerchoice}
  Acceleration magnitude in m/s² printed beside each displayed choice
\end{definition}
\begin{proof}
  This is the declared model definition of acceleration Meters Per Second Squared.
\end{proof}
\begin{definition}[recorded Answer Choice]
  \label{def:physics:phyx-mini-0755:phyxminiproblems-problemphyxmini0755-recordedanswerchoice}
  \lean{PhyXMiniProblems.ProblemPhyXMini0755.recordedAnswerChoice}
  \uses{def:physics:phyx-mini-0755:phyxminiproblems-problemphyxmini0755-answerchoice}
  The answer label recorded in the supplied dataset metadata
\end{definition}
\begin{proof}
  This is the declared model definition of recorded Answer Choice.
\end{proof}
\begin{definition}[answer Choice Error Meters Per Second Squared]
  \label{def:physics:phyx-mini-0755:phyxminiproblems-problemphyxmini0755-answerchoiceerrormeterspersecondsquared}
  \lean{PhyXMiniProblems.ProblemPhyXMini0755.answerChoiceErrorMetersPerSecondSquared}
  \uses{def:physics:phyx-mini-0755:phyxminiproblems-problemphyxmini0755-accelerationmagnitudeinmeterspersecondsquared, def:physics:phyx-mini-0755:phyxminiproblems-problemphyxmini0755-asteroidpushsetup, def:physics:phyx-mini-0755:phyxminiproblems-problemphyxmini0755-answerchoice}
  Absolute discrepancy from one displayed acceleration magnitude
\end{definition}
\begin{proof}
  This is the declared model definition of answer Choice Error Meters Per Second Squared.
\end{proof}
\begin{definition}[Matches Displayed Acceleration]
  \label{def:physics:phyx-mini-0755:phyxminiproblems-problemphyxmini0755-matchesdisplayedacceleration}
  \lean{PhyXMiniProblems.ProblemPhyXMini0755.MatchesDisplayedAcceleration}
  \uses{def:physics:phyx-mini-0755:phyxminiproblems-problemphyxmini0755-asteroidpushsetup, def:physics:phyx-mini-0755:phyxminiproblems-problemphyxmini0755-answerchoice, def:physics:phyx-mini-0755:phyxminiproblems-problemphyxmini0755-answerchoiceerrormeterspersecondsquared}
  Agreement with a value displayed to the nearest 0.01 m/s²
\end{definition}
\begin{proof}
  This is the declared model definition of Matches Displayed Acceleration.
\end{proof}
\begin{definition}[Is Unique Closest Acceleration Choice]
  \label{def:physics:phyx-mini-0755:phyxminiproblems-problemphyxmini0755-isuniqueclosestaccelerationchoice}
  \lean{PhyXMiniProblems.ProblemPhyXMini0755.IsUniqueClosestAccelerationChoice}
  \uses{def:physics:phyx-mini-0755:phyxminiproblems-problemphyxmini0755-asteroidpushsetup, def:physics:phyx-mini-0755:phyxminiproblems-problemphyxmini0755-answerchoice, def:physics:phyx-mini-0755:phyxminiproblems-problemphyxmini0755-answerchoiceerrormeterspersecondsquared}
  A displayed value is strictly closer than every alternative
\end{definition}
\begin{proof}
  This is the declared model definition of Is Unique Closest Acceleration Choice.
\end{proof}
\begin{lemma}[acceleration Magnitude eq exact Expression]
  \label{lem:physics:phyx-mini-0755:phyxminiproblems-problemphyxmini0755-accelerationmagnitude-eq-exactexpression}
  \lean{PhyXMiniProblems.ProblemPhyXMini0755.accelerationMagnitude_eq_exactExpression}
  \uses{def:physics:phyx-mini-0755:phyxminiproblems-problemphyxmini0755-accelerationmagnitudeinmeterspersecondsquared, def:physics:phyx-mini-0755:phyxminiproblems-problemphyxmini0755-asteroidpushsetup, def:physics:phyx-mini-0755:phyxminiproblems-problemphyxmini0755-matchesasteroidguidancescenario, def:physics:phyx-mini-0755:phyxminiproblems-problemphyxmini0755-matchesproblemreadouts, def:physics:phyx-mini-0755:phyxminiproblems-problemphyxmini0755-matchesprimaryfigure, def:physics:phyx-mini-0755:phyxminiproblems-problemphyxmini0755-matchesfigureforcegeometry, def:physics:phyx-mini-0755:phyxminiproblems-problemphyxmini0755-satisfiesnewtoniandynamics, def:physics:phyx-mini-0755:phyxminiproblems-problemphyxmini0755-exactaccelerationmagnitudeinmeterspersecondsquared}
  The governing assumptions determine the exact unrounded acceleration
\end{lemma}
\begin{proof}
  The conclusion follows from the governing relations and figure readouts named in the statement of acceleration Magnitude eq exact Expression.
\end{proof}
\begin{lemma}[exact Acceleration selects choice C]
  \label{lem:physics:phyx-mini-0755:phyxminiproblems-problemphyxmini0755-exactacceleration-selects-choice-c}
  \lean{PhyXMiniProblems.ProblemPhyXMini0755.exactAcceleration_selects_choice_C}
  \uses{def:physics:phyx-mini-0755:phyxminiproblems-problemphyxmini0755-accelerationmagnitudeinmeterspersecondsquared, def:physics:phyx-mini-0755:phyxminiproblems-problemphyxmini0755-asteroidpushsetup, def:physics:phyx-mini-0755:phyxminiproblems-problemphyxmini0755-exactaccelerationmagnitudeinmeterspersecondsquared, def:physics:phyx-mini-0755:phyxminiproblems-problemphyxmini0755-matchesdisplayedacceleration, def:physics:phyx-mini-0755:phyxminiproblems-problemphyxmini0755-isuniqueclosestaccelerationchoice}
  The exact value is approximately 0.8753 m/s², so it rounds to the displayed value 0.88 m/s² and that value is uniquely closest among the four choices
\end{lemma}
\begin{proof}
  The conclusion follows from the governing relations and figure readouts named in the statement of exact Acceleration selects choice C.
\end{proof}
% --- Archon declaration topology end ---
% --- Archon physics formalization source end ---
